\chapter{جمع‌بندی و کارهای آتی}\label{Chap:Chap5}
\minitoc
\setcounter{footnote}{0}
\pagestyle{plain}

در فصل 
\ref{Chap:Chap1}
به معرفی مسئله‌ی دسته‌بندی بزرگ مقیاس نمونه محدود و اهمیت آن پرداختیم و به رویکرد حل‌مسئله و چالش‌های موجود در این حوزه  اشاره کردیم. در فصل 
 \ref{Chap:Chap2}
رویکرد‌ها و چالش‌های اصلی این تنظیمات یادگیری را در قالب مرور مقالات مهم هر حوزه بیان کردیم. در واقع در این فصل با مرور پژوهش‌های پیشین، به نوعی ادبیات و دانش پایه مرتبط با این حوزه را در کنار بیان عملکرد هر کدام از مقالات معرفی شد.

در فصل
\ref{Chap:Chap3}
راهکار پیشنهادی را به طور کامل توضیح دادیم. به این منظور مدل‌سازی تئوری که از مسئله داشتیم در قالب روابط ریاضی بیان شد و برای واضح شدن رویکرد یادگیری، شبه‌کدهای مرتبط ارائه شد. در نهایت در فصل 
\ref{Chap:Chap4}
به بررسی کیفیت روش پیشنهادی به صورت کمی و در مقایسه با سایر روش‌ها پرداختیم. از طرفی معماری دقیق مدل‌های دسته‌بند درشت‌دانه و ریز‌دانه در کنار مجموعه دادگانی که مورد استفاده قرار گرفتند، معرفی شدند. در این فصل، به جمع‌بندی موارد گفته‌شده می‌پردازیم و چالش‌های باقی مانده و حوزه‌هایی که هنوز جای کار دارند را معرفی می‌کنیم.

\section{جمع‌بندی و نقاط قوت روش}

در این قسمت مزیت‌های کلیدی رویکرد پیشنهادی در حل چالش‌های مسئله‌ی یادگیری بزرگ مقیاس نمونه محدود بیان شده‌اند.

\begin{itemize}
	\item 
	با ارائه‌ی یک تنظیمات جدید در مسئله‌ی یادگیری بزرگ مقیاس نمونه محدود، سخت‌ترین ستینگ ممکن را انتخاب کردیم که با ناسازگار بودن داده‌ها و برچسب‌های مجموعه دادگان مبدا و مقصد تشکیل می‌شود. در واقع این تنظیماتی است که به واقعیت نزدیک‌تر است چرا که چالش‌های مجموعه باز بودن و تغییر توزیع زیرجمعیت را دربرمی‌گیرد.
	
	\item 
	در سطح درشت دانه، یک روش بدیع برای مقاومت در برابر تغییر توزیع زیرجمعیت به صورت نیمه‌نظارتی پیشنهاد دادیم که به آن تنظیم ترکیبی می‌گوییم.
	
	\item 
	
	از بدنه‌های بنیادی بیانگر بهره بردیم و به همین جهت لایه‌هایی که آموزش می‌دهیم تا حد ممکن کوچک اند. در نتیجه هم بهینگی محاسباتی را سبب می‌شوند و هم در تنظیمات خاص این مسئله که تغییر توزیع و نمونه‌ی محدود داریم، از بیش برازش روی مجموعه‌ی مبدا جلوگیری می‌شود.
	
	\item 
	از سلسله‌مراتب ذاتی در داده‌‌ها استفاده کردیم و در یک چارچوب بیزی، به ترکیب اطلاعات درشت‌دانه برای دسته‌بندی در سطح ریزدانه پرداختیم. به این ترتیب یک سوگیری القایی مبتنی بر وجود سلسله‌مراتب در یادگیری داریم که در بازنمایی داده‌ها بروز پیدا می‌کند و باعث منظم شدن فضای دسته‌ها در شرایط مجموعه باز می ‌شود.
	
	\item 
با توجه به مشاهدات، از یک راهبرد یکی از همه برنده استفاده کردیم که تعداد دسته‌های هدف را محدود میکند و باعث بالا رفتن دقت می‌شود.
	
	\item 
با تعریف ساختار سلسله‌مراتبی، سبب شدیم در سطح درشت‌دانه با وجود نمونه محدود بودن دسته‌ها، داده‌ی کافی برای آموزش به صورت تمایزگذار فراهم شود و تمامی روش های متداول دسته‌بندی در یادگیری ژرف را بتوان در این سطح استفاده کرد.

	\item
در مجموعه‌ی مبدا هم علاوه بر مقصد، دسته‌ها تعداد محدودی نمونه برچسب خورده دارند. برخلاف اکثر روش‌های دیگر که نیاز به یک مجموعه‌ی داده‌ی مبدا با تعداد بالای داده‌ی آموزشی دارند.
	
	\item 
در سطح ریزدانه، از یک رویکرد متایادگیری مولد استفاده کردیم که در عین داشتن مزایای متایادگیری، قابلیت کار در شرایط مجموعه باز و ظهور دسته‌های بدیع را دارد.
	
	\item 
رویکرد دسته‌بندی ما در سطح ریزدانه، یک یادگیری
\trans{چند وظیفه‌ای}{Multi Task}
 هست و در واقع چیزی که مدل بر روی آن مشروط می‌شود، یک شناسه‌ی وظیفه است. بنابراین سر‌های قابل آموزش، از دانش مشترک بین ابردسته‌ها یاد می‌گیرند و  یک رویکرد 
\trans{ترکیب خبره‌ها }{Mixture of Experts}
داریم.
	
	\item 
	با توجه به نتایج ارائه شده، رویکرد پیشنهادی در قالب دسته‌بندی سلسله‌مراتبی می‌تواند در تنظیمات
	\trans{بدون یادگیری}{Free Training}
	 هم به کارایی لازم برسد. به این منظور لازم است از بدنه‌های بنیادی مناسبی استفاده شود.
	
\end{itemize}

\section{کار‌های آتی}

همانطور که در جداول فصل
\ref{Chap:Chap4}
دیدیم، رویکرد پیشنهادی بر روی بدنه‌های کوچک و متوسط منجر به افزایش دقت بسیار زیادی می‌شود. همچنین روی بدنه‌های بنیادی بزرگ مانند مدل CLIP که با حجم بسیار بالای داده آموزش دیده، به افزایش دقت قابل توجهی در تمامی شرایط تعریف شده اعم از با یادگیری و بدون یادگیری می‌رسیم. با این حال در بدنه‌های بنیادی بسیار قوی مانند Dino که از طرف شرکت‌های صنعتی بزرگ با تعداد پارامترهای بسیار بالا و آموزش روی مجموعه داده‌های بزرگ ارائه شده‌اند، افزایش دقت وجود دارد اما به اندازه‌ی موارد بیان شده‌ی دیگر نیست. با این وجود، اگر از دقت حاصل از مدل 
\trans{پیشگو}{Oracle}
در سطح درشت‌دانه استفاده کنیم، دیدیم که در مدل‌های بنیادی خیلی بزرگ هم می‌توان افزایش دقت خوبی داشت. با توجه به این که دسته‌بند سطح درشت‌دانه‌ی ما، تنها یک لایه‌ی خطی بر روی بدنه‌ی پیش‌آموزش دیده آموزش می‌دهد، به نظر می‌رسد با کار بیشتر روی دسته‌بند این سطح بتوان به دقت حدود دسته‌بند پیشگو رسید. این موضوع با وجود تجمیع نمونه‌های دسته‌های پایه در سطح درشت دانه و تعداد کم دسته‌ها در حدود زیر ۱۰۰ ابردسته، قابل دستیابی است.

مورد دیگری که جای بهبود دارد، در نظر گرفتن اضافه شدن ابردسته‌های جدید در بعضی تنظیمات خیلی خاص است. به این منظور می‌توان روش‌های حکاکی وزن‌ها که در قسمت
\ref{}
اشاره شد را مورد توجه قرار داد. در این صورت همچنان می‌توان از رویکرد یادگیری تمایزگذار و اکثر روش‌های متداول یادگیری ژرف مستقل از معماری و شرایط دیگر بهره‌برد.
